\section{The cluster state}\label{sec:cluster}
%% ===================================================================

The cluster state~\cite{BriegelRaussendorf2001,RaussendorfBriegel2001} is a
universal resource for measurement-based quantum computation; its
one-dimensional chain provides another canonical example of a
non-trivial SPT phase~\cite{Schuch2011Phases,Cirac2021Matrix}. We follow
the presentation of~\cite[Appendix~A, ``The cluster state'']{Cirac2021Matrix}.
Throughout, $\ket{\pm} = \tfrac{1}{\sqrt{2}}(\ket{0} \pm \ket{1})$.

\begin{definition}[Cluster-state tensor of the review]\label{def:cluster_tensor_review}
    \lean{MPSTensor.clusterTensorRMP}
    \leanok
    \uses{def:mps_tensor}
    Set $d = D = 2$. The cluster-state tensor
    of~\cite[Appendix~A, ``The cluster state'']{Cirac2021Matrix} is the family
    $\{A^0, A^1\} \subset \MN{2}$ given by
    \begin{align}
        A^0
         & = \ketbra{0}{+}
        = \frac{1}{\sqrt{2}}
        \begin{pmatrix} 1 & 1 \\ 0 & 0 \end{pmatrix},
        \qquad
        A^1
        = \ketbra{1}{-}
        = \frac{1}{\sqrt{2}}
        \begin{pmatrix} 0 & 0 \\ 1 & -1 \end{pmatrix}.
        \notag
    \end{align}
    Equivalently, $(A^s)_{\alpha\beta} = \delta_{\alpha s}\,(-1)^{s\beta}/\sqrt{2}$.
\end{definition}

\begin{theorem}[Controlled-$Z$ construction of the cluster state]\label{thm:cluster_cz_construction}
    \lean{MPSTensor.clusterTensorRMP_mpv_eq_clusterCZRing,
        MPSTensor.clusterCZRing, MPSTensor.plusProductState,
        MPSTensor.controlledZ}
    \leanok
    \uses{def:cluster_tensor_review}
    For every $N \geq 1$ and every $s \in \{0,1\}^N$,
    \begin{align}
        \tr\bigl(A^{s_1} A^{s_2} \cdots A^{s_N}\bigr)
        = \Bigl(\,\prod_{j=1}^{N} \mathrm{CZ}_{j,j+1}\,
        \ket{+}^{\otimes N}\Bigr)_{s}
        = 2^{-N/2} \prod_{j=1}^{N} (-1)^{s_j s_{j+1}},
        \notag
    \end{align}
    with indices modulo $N$, where $\mathrm{CZ}_{a,b}$ multiplies the
    coefficient of $\ket{s}$ by $(-1)^{s_a s_b}$. Thus the periodic matrix
    product vector of the tensor of
    Definition~\ref{def:cluster_tensor_review} is exactly the normalized
    cluster state obtained by acting with controlled-$Z$ gates between nearest
    neighbours of a ring on
    $\ket{+}^{\otimes N}$~\cite[Appendix~A, ``The cluster state'']{Cirac2021Matrix},
    \cite{RaussendorfBriegel2001}. For $N \geq 3$ the $N$ gates act on the
    distinct bonds of the ring. For $N = 2$ the two gates act on the same pair
    and cancel, giving $\ket{+}^{\otimes 2}$; for $N = 1$ the single gate
    $\mathrm{CZ}_{1,1}$ is the Pauli $Z$, giving $\ket{-}$.
\end{theorem}

\begin{proof}
    \leanok
    Since $A^s = \ketbra{s}{h_s}$ with $h_0 = +$ and $h_1 = -$, the trace of
    the product is $\prod_j \braket{h_{s_j}}{s_{j+1}}$ with indices modulo
    $N$, and $\braket{h_a}{b} = (-1)^{ab}/\sqrt{2}$. On the other side,
    every coefficient of $\ket{+}^{\otimes N}$ equals $2^{-N/2}$, and the
    gates multiply the coefficient of $\ket{s}$ by
    $\prod_j (-1)^{s_j s_{j+1}}$.
\end{proof}

\begin{lemma}[Valence-bond form of the cluster-state tensor]\label{lem:cluster_valence_bond}
    \lean{MPSTensor.clusterTensorVBS, MPSTensor.clusterTensorVBS_mpv_eq_prod,
        MPSTensor.clusterTensorVBS_eq_smul}
    \leanok
    \uses{def:cluster_tensor_review}
    In the valence-bond description of~\cite{VerstraeteCirac2004ValenceBond},
    neighbouring virtual qubits share the bond
    $\ket{00} + \ket{01} + \ket{10} - \ket{11}$ and each site maps both of its
    virtual qubits to the physical one by
    $\ket{0}\bra{00} + \ket{1}\bra{11}$, so that the coefficient of $\ket{s}$
    on a ring of $N \geq 1$ sites is $\prod_{j=1}^{N} H_{s_j s_{j+1}}$ with
    $H = \bigl(\begin{smallmatrix} 1 & 1 \\ 1 & -1 \end{smallmatrix}\bigr)$
    and indices modulo $N$. Let $B^s = \ketbra{s}{s} H$. Then for every
    $N \geq 1$ and every $s \in \{0,1\}^N$,
    \begin{align}
        \tr\bigl(B^{s_1} B^{s_2} \cdots B^{s_N}\bigr)
        = \prod_{j=1}^{N} H_{s_j s_{j+1}},
        \notag
    \end{align}
    so $B$ is a matrix product description of the valence-bond state, and
    $B^s = \sqrt{2}\,A^s$.
\end{lemma}

\begin{proof}
    \leanok
    The entries are $(B^s)_{\alpha\beta} = \delta_{\alpha s} H_{s\beta}$, so in
    the expansion of the trace as a sum over cyclic index paths only the path
    $\alpha_j = s_j$ survives, and it contributes
    $\prod_j H_{s_j s_{j+1}}$. The identity $B^s = \sqrt{2}\,A^s$ is read off
    from $\ketbra{s}{s}H = \sqrt{2}\,\ketbra{s}{h_s}$ with $h_0 = +$ and
    $h_1 = -$.
\end{proof}

\begin{definition}[Transposed representative]\label{def:cluster_tensor}
    \lean{MPSTensor.clusterTensor}
    \leanok
    \uses{def:mps_tensor}
    The results below are proved for the transpose of the tensor of
    Definition~\ref{def:cluster_tensor_review}, the family
    \begin{align}
        A^0
         & = \ketbra{+}{0}
        = \frac{1}{\sqrt{2}}
        \begin{pmatrix} 1 & 0 \\ 1 & 0 \end{pmatrix},
        \qquad
        A^1
        = \ketbra{-}{1}
        = \frac{1}{\sqrt{2}}
        \begin{pmatrix} 0 & 1 \\ 0 & -1 \end{pmatrix},
        \notag
    \end{align}
    and transferred back by Lemma~\ref{lem:cluster_review_gauge}.
\end{definition}

\begin{lemma}[The two cluster-state tensors are gauge equivalent]\label{lem:cluster_review_gauge}
    \lean{MPSTensor.clusterTensorRMP_eq_transpose,
        MPSTensor.clusterTensor_gaugeEquiv_clusterTensorRMP,
        MPSTensor.clusterTensor_sameMPV_clusterTensorRMP,
        MPSTensor.clusterBlocked_gaugeEquiv_clusterBlockedRMP}
    \leanok
    \uses{def:cluster_tensor_review, def:cluster_tensor, def:gauge_equiv}
    The tensor of Definition~\ref{def:cluster_tensor_review} is the transpose
    of the tensor of Definition~\ref{def:cluster_tensor}, and the two are
    gauge equivalent: with
    $H = \bigl(\begin{smallmatrix} 1 & 1 \\ 1 & -1 \end{smallmatrix}\bigr)$,
    the review tensor is $H A^s H^{-1}$, since $H\ket{\pm} = \sqrt{2}\,\ket{0}$
    and $\sqrt{2}\,\ket{1}$ respectively. Hence both tensors generate the same
    matrix product vectors at every length, and their length-$2$ blockings are
    gauge equivalent by the same matrix.
\end{lemma}

\begin{theorem}[Properties of the review cluster-state tensor]\label{thm:cluster_review_transferred}
    \lean{MPSTensor.clusterTensorRMP_not_isInjective,
        MPSTensor.clusterTensorRMP_isNBlkInjective_two,
        MPSTensor.clusterTensorRMP_isNormal,
        MPSTensor.clusterBlockedRMP, MPSTensor.clusterBlockedRMP_isInjective,
        MPSTensor.clusterBlockedRMP_isOnSiteSymmetric_Z2Z2,
        MPSTensor.clusterRepXRMP, MPSTensor.clusterProjRepRMP,
        MPSTensor.clusterBlockedRMP_twist_intertwine,
        MPSTensor.clusterBlockedRMP_hasStringOrder,
        MPSTensor.clusterBlockedRMP_isZCL}
    \leanok
    \uses{lem:cluster_review_gauge, thm:cluster_not_injective, thm:cluster_normal,
        thm:cluster_z2z2_symmetric, def:cluster_omega, thm:cluster_nontrivial_spt,
        thm:cluster_string_order, thm:cluster_zcl}
    The tensor of Definition~\ref{def:cluster_tensor_review} is not injective,
    its length-$2$ blocking is injective, and so it is normal. The length-$2$
    blocked tensor is on-site symmetric under the $\Z_2 \times \Z_2$ action of
    Theorem~\ref{thm:cluster_z2z2_symmetric}; the two generators are
    implemented on the bond space by $\sigma_x$ and $\sigma_z$, the
    conjugates by $H$ of the gauges of
    Theorem~\ref{thm:cluster_gauge_anticomm}, and these form a projective
    representation with the factor system of
    Definition~\ref{def:cluster_omega}, whose class is non-trivial
    (Theorem~\ref{thm:cluster_nontrivial_spt}). The blocked tensor has string
    order for every group element with the maximally mixed boundary state, and
    has zero correlation length.
\end{theorem}

\begin{proof}
    \leanok
    Each property is invariant under the gauge equivalence of
    Lemma~\ref{lem:cluster_review_gauge} and holds for the transposed
    representative by Theorems~\ref{thm:cluster_not_injective},
    \ref{thm:cluster_normal}, \ref{thm:cluster_z2z2_symmetric},
    \ref{thm:cluster_string_order} and~\ref{thm:cluster_zcl}. The twisted
    blocked tensors and their virtual action are the conjugates by $H$ of
    those of the transposed representative, so the intertwining relations
    carry over. For the string order, the blocked transfer channel of the
    transposed representative is unital and its adjoint fixes the maximally
    mixed state; since $H$ is real symmetric with $H^2 = 2\,\Id$, conjugation
    by $H$ carries both properties to the blocked tensor of
    Definition~\ref{def:cluster_tensor_review}.
\end{proof}

\begin{theorem}[Cluster-state tensor is not injective]\label{thm:cluster_not_injective}
    \lean{MPSTensor.cluster_not_isInjective}
    \leanok
    \uses{def:cluster_tensor, def:injective}
    The cluster-state tensor is not injective:
    $\spn_{\C}\{A^0, A^1\} = \spn_{\C}\{\ketbra{+}{0}, \ketbra{-}{1}\}$
    has dimension~$2$, not $\dim \MN{2} = 4$.
\end{theorem}

\begin{proof}
    \leanok
    \uses{def:cluster_tensor}
    Both $A^0$ and $A^1$ are rank-one matrices with linearly
    independent column spaces, so they are linearly independent.
    Their span is two-dimensional, which is strictly less than
    $\dim \MN{2} = 4$.
\end{proof}

\begin{theorem}[Cluster-state tensor is normal]\label{thm:cluster_normal}
    \lean{MPSTensor.cluster_isNBlkInjective_two,
        MPSTensor.clusterBlocked_isInjective}
    \leanok
    \uses{def:cluster_tensor, def:injective}
    Although the cluster-state tensor is not injective at a single site, its
    length-$2$ blocking is injective: the four length-$2$ products span
    $\MN{2}$. In particular, the cluster-state tensor is normal.
\end{theorem}

\begin{remark}\label{rem:cluster_injective_blocked}
    Although the cluster-state tensor is not injective at length~$1$,
    the length-$2$ blocked tensor is injective. Using
    $\braket{0}{\pm} = \tfrac{1}{\sqrt{2}}$,
    $\braket{1}{+} = \tfrac{1}{\sqrt{2}}$, and
    $\braket{1}{-} = -\tfrac{1}{\sqrt{2}}$, the
    four products $A^{ij} = A^i A^j$ are
    \begin{align}
        A^{00} & = \tfrac{1}{\sqrt{2}}\ketbra{+}{0},\quad
        A^{01} = \tfrac{1}{\sqrt{2}}\ketbra{+}{1},\quad
        A^{10} = \tfrac{1}{\sqrt{2}}\ketbra{-}{0},\quad
        A^{11} = -\tfrac{1}{\sqrt{2}}\ketbra{-}{1}.
        \label{eq:ex_cluster_blocked_products}
    \end{align}
    Since $\{\ket{+},\ket{-}\}$ and $\{\bra{0},\bra{1}\}$ are bases,
    these four rank-one matrices are linearly independent and
    span~$\MN{2}$.
\end{remark}

\begin{theorem}[$\Z_2 \times \Z_2$ on-site symmetry of the blocked cluster-state tensor]\label{thm:cluster_z2z2_symmetric}
    \lean{MPSTensor.cluster_isOnSiteSymmetric_Z2Z2}
    \leanok
    \uses{def:cluster_tensor, def:on_site_symmetric}
    The length-$2$ blocked cluster-state tensor is on-site symmetric
    under $\Z_2 \times \Z_2$. The two generators act on the
    $4$-dimensional blocked physical space $(\C^2)^{\otimes 2}$ by
    $\sigma_x \otimes \Id$ and $\Id \otimes \sigma_x$,
    which commute and each square to the identity, defining a linear
    representation of $\Z_2 \times \Z_2$.
\end{theorem}

\begin{theorem}[Anticommuting virtual gauges of the cluster-state symmetry]\label{thm:cluster_gauge_anticomm}
    \lean{MPSTensor.cluster_gauge_anticomm}
    \leanok
    The virtual matrices implementing the two generators of
    Theorem~\ref{thm:cluster_z2z2_symmetric}, namely $V_1 = \sigma_z$ and
    $V_2 = \sigma_x$, anticommute ($V_1 V_2 = -V_2 V_1$).
\end{theorem}

\begin{definition}[Cluster-state factor system]\label{def:cluster_omega}
    \lean{MPSTensor.clusterOmega}
    \leanok
    \uses{def:cluster_tensor}
    The cluster-state factor system on $\Z_2 \times \Z_2$ is
    $\omega(g,h) = (-1)^{g_2\,h_1}$ (the value $-1$ when $g_2 = 1$ and
    $h_1 = 1$, and $1$ otherwise). It is the factor system of the explicit
    projective representation with virtual gauges $\sigma_z$ and $\sigma_x$,
    both squaring to the identity, so there is no diagonal term.
\end{definition}

\begin{lemma}[Cluster-state commutator phase]\label{lem:cluster_comm_phase_neg_one}
    \lean{MPSTensor.cluster_commPhase_eq_neg_one}
    \leanok
    \uses{def:cluster_omega, def:comm_phase}
    The commutator phase of the cluster-state factor system on the two generators
    is $-1$.
\end{lemma}

\begin{theorem}[The cluster-state factor system is a $2$-cocycle]\label{thm:cluster_omega_cocycle}
    \lean{MPSTensor.clusterOmega_isCocycle}
    \leanok
    \uses{def:is_cocycle, def:cluster_omega}
    The cluster-state factor system is a genuine $2$-cocycle, being the factor
    system of an explicit projective representation, so it represents an element
    of $H^2(\Z_2 \times \Z_2, \mathrm{U}(1))$.
\end{theorem}

\begin{theorem}[The cluster-state $\Z_2 \times \Z_2$ phase is non-trivial]\label{thm:cluster_nontrivial_spt}
    \lean{MPSTensor.cluster_isNontrivialSPT}
    \leanok
    \uses{def:nontrivial_class, thm:nontrivial_of_comm_phase, lem:cluster_comm_phase_neg_one}
    The cluster-state factor system (Theorem~\ref{thm:cluster_omega_cocycle}),
    built from the anticommuting gauges $\sigma_z$ and $\sigma_x$
    (Theorem~\ref{thm:cluster_gauge_anticomm}), is not a coboundary: its
    commutator phase on the two generators is $-1$. Its class in
    $H^2(\Z_2 \times \Z_2, \mathrm{U}(1)) \cong \Z_2$ is therefore non-trivial, so
    the cluster state realizes a non-trivial SPT phase.
\end{theorem}

\begin{theorem}[The cluster state has string order under its $\Z_2 \times \Z_2$ symmetry]\label{thm:cluster_string_order}
    \lean{MPSTensor.cluster_hasStringOrder}
    \leanok
    \uses{def:has_string_order, def:cluster_tensor}
    For every element $g$ of $\Z_2 \times \Z_2$, the length-$2$ blocked
    cluster-state tensor has string order in the sense of
    Definition~\ref{def:has_string_order} for the twist by $g$, with the
    maximally mixed boundary state $\Lambda = \tfrac{1}{2}\Id$ as its stationary
    boundary state.
\end{theorem}
\begin{proof}
    \leanok
    \uses{thm:has_string_order_of_symmetric_injective, thm:cluster_z2z2_symmetric, rem:cluster_injective_blocked}
    The blocked tensor is injective (Remark~\ref{rem:cluster_injective_blocked})
    and on-site symmetric under $\Z_2 \times \Z_2$
    (Theorem~\ref{thm:cluster_z2z2_symmetric}), and each group element acts by a
    real symmetric involutive permutation matrix, hence unitarily. The
    maximally mixed state $\Lambda = \tfrac{1}{2}\Id$ is positive definite and
    has trace~$1$. The four blocked matrices in
    (\ref{eq:ex_cluster_blocked_products}) satisfy
    $\sum_i A^i (A^i)^\dagger = \Id$ and
    $\sum_i (A^i)^\dagger A^i = \Id$, so the transfer map is unital and
    $\Lambda$ is a fixed point of the adjoint transfer map. Applying
    Theorem~\ref{thm:has_string_order_of_symmetric_injective} yields string
    order for every $g$.
\end{proof}

\begin{theorem}[Blocked cluster transfer map is idempotent]\label{thm:cluster_blocked_idempotent}
    \lean{MPSTensor.clusterBlocked_transferMap_idempotent,
        MPSTensor.clusterBlocked_transferMap_apply}
    \leanok
    \uses{def:cluster_tensor, def:transfer_map}
    The length-$2$ blocked cluster transfer map satisfies $\E_A^2 = \E_A$. In
    closed form, it sends every $X$ to the scalar $(X_{00} + X_{11})/2$ times
    the identity, so its image is one-dimensional.
\end{theorem}

\begin{proof}
    \leanok
    The single-site cluster transfer map is not idempotent, but the four blocked
    matrices give the closed form
    $\E_A(X) = \tfrac{1}{2}(X_{00} + X_{11}) \Id$, which is manifestly
    idempotent.
\end{proof}

\begin{theorem}[Cluster zero correlation length]\label{thm:cluster_zcl}
    \lean{MPSTensor.clusterBlocked_isZCL}
    \leanok
    \uses{def:is_zcl, thm:cluster_blocked_idempotent, thm:zcl_iff_idempotent_transfer}
    The length-$2$ blocked cluster tensor has zero correlation length: its
    idempotent transfer map (Theorem~\ref{thm:cluster_blocked_idempotent}) gives
    correlations independent of the separation. Like the GHZ state, and unlike
    the AKLT state, the cluster state is a renormalization fixed point with zero
    correlation length.
\end{theorem}

\begin{proof}
    \leanok
    \uses{thm:zcl_iff_idempotent_transfer, thm:cluster_blocked_idempotent}
    The claim follows from the idempotence of the blocked cluster transfer map
    and the equivalence of zero correlation length with transfer-map idempotence
    (Theorem~\ref{thm:zcl_iff_idempotent_transfer}).
\end{proof}

%% ===================================================================
